\chapter{Blocks 3--4: Python for Genomic Data}

\textbf{Primary:} Data structures, Pandas, visualization\quad|\quad\textbf{Interleave:} Unix pipes, real genomic files

\begin{brainbox}[Connection to what you already know]
You just learned to process text files with grep and awk. Python does the same things but with more power, reusable code, and the ability to make plots and run statistics.
\end{brainbox}


\section{Block 3: Core Python via Biology}

\sprint{1}{Data Structures Are Biology}{25}
Pythonic code is less about clever tricks and more about reaching for the right built-in collection. Each has a natural biological use case.
\begin{itemize}[nosep]
  \item \textbf{Strings = DNA sequences.}
    \begin{lstlisting}
seq = 'ATGCGATCG'
comp = seq.translate(str.maketrans('ATCG','TAGC'))
print(comp[::-1])  # reverse complement
    \end{lstlisting}
  \item \textbf{Lists = ordered, mutable collections.} Comprehensions are how you filter them: \texttt{[g for g in genes if g.startswith('BRCA')]}. Swap the brackets for parentheses to get a \textit{generator expression} that yields items one at a time without building the full list in memory.
  \item \textbf{Tuples = immutable records.} Genomic coordinates are a natural fit: \texttt{peak = ("chr1", 1000, 2000)}. Because tuples are immutable, they're hashable --- so they can serve as dict keys or live in sets (``which regions have I already processed?'').
  \item \textbf{Dictionaries = lookup tables.} Build a codon table: \texttt{codons = \{'ATG':\,'M', 'TAA':\,'*', 'GCT':\,'A', \ldots\}} and translate three bases at a time.
  \item \textbf{Sets = membership and overlap.} \texttt{set(sampleA) \& set(sampleB)} (intersection), \texttt{|} (union), \texttt{-} (difference). Use them when order doesn't matter and you care about presence/absence.
  \item \textbf{\texttt{collections.Counter} and \texttt{defaultdict}.} \texttt{Counter(seq).most\_common(4)} counts every character and returns the top four --- no ``does this key exist yet?'' dance. \texttt{defaultdict(list)} lets you write \texttt{d[chrom].append(pos)} without initializing each chromosome's list by hand.
\end{itemize}
\begin{winbox}[Checkpoint]
You reverse-complemented DNA with \texttt{str.translate}, counted nucleotide frequencies with \texttt{Counter}, and grouped positions by chromosome with \texttt{defaultdict}. Idiomatic Python is about picking the right container, not writing loops.
\end{winbox}

\sprint{2}{Parsing Genomic Files}{25}
\begin{itemize}[nosep]
  \item Write a FASTA parser: open the file, loop through lines, yield (header, sequence) pairs
  \item Use a context manager (\texttt{with open(...) as f:}) --- always. This prevents resource leaks.
  \item Parse a BED file into a list of dictionaries: \texttt{[\{chr: ..., start: ..., end: ...\}, ...]}
  \item Handle errors: what happens when a line is malformed? Add \texttt{try/except}.
\end{itemize}
\begin{winbox}[Checkpoint]
You can parse any text-based genomic file. This is a skill you'll use indefinitely.
\end{winbox}

\sprint{3}{Generators --- Process Files Bigger Than Your RAM}{25}
\begin{itemize}[nosep]
  \item A generator uses \texttt{yield} instead of \texttt{return} --- it produces items one at a time without loading everything into memory
  \item Convert your FASTA parser to a generator. Notice: it handles a 10GB file the same way as a 10KB file.
  \item This is critical for bioinformatics: VCFs can be billions of lines. You cannot load them into memory.
\end{itemize}
\begin{winbox}[Checkpoint]
Your FASTA parser works on any size file. That's production-quality code.
\end{winbox}

\sprint{4}{Pythonic Idioms \& Type Hints}{25}
Writing Python that other people (and future-you) can read quickly is less about clever tricks and more about using the idioms the language was built around. Each of these takes five minutes to learn and saves hours of debugging.
\begin{itemize}[nosep]
  \item \textbf{f-strings.} \texttt{f"\{gene\} on chr\{chrom\} (padj=\{pval:.2e\})"} --- variables, formatting, and expressions in one readable line. Use these exclusively for user-facing strings; they're faster than \texttt{\%} or \texttt{.format()} and much harder to misread.
  \item \textbf{\texttt{pathlib} for filesystem paths.}
    \begin{lstlisting}
from pathlib import Path
data = Path.home() / "bioinfo" / "block1" / "data"
for vcf in data.glob("*.vcf"):
    print(vcf.stem, vcf.stat().st_size)
    \end{lstlisting}
    Path objects beat string concatenation: they handle OS differences, support \texttt{/} for joining, and expose \texttt{.stem}, \texttt{.suffix}, \texttt{.parent}, \texttt{.exists()}.
  \item \textbf{\texttt{enumerate} and \texttt{zip}.} Almost never write \texttt{for i in range(len(xs))} --- use \texttt{for i, x in enumerate(xs)} or \texttt{for a, b in zip(xs, ys)}. These are the canonical ways to walk sequences.
  \item \textbf{Dataclasses for records.}
    \begin{lstlisting}
from dataclasses import dataclass

@dataclass
class Variant:
    chrom: str
    pos: int
    ref: str
    alt: str
    qual: float = 0.0
    \end{lstlisting}
    You get \texttt{\_\_init\_\_}, \texttt{\_\_repr\_\_}, and equality for free, with attribute access (\texttt{v.pos}) instead of \texttt{v["pos"]}. Prefer these over parallel lists or dicts-of-fields for anything with a fixed shape.
  \item \textbf{Type hints.} Annotations like \texttt{def parse\_vcf(path: Path) -> list[Variant]:} are not enforced at runtime, but tools (\texttt{mypy}, Pyright, your editor) use them to catch bugs before you run the code, and they're the best function documentation you can write. Start with hints on function signatures and dataclass fields --- don't try to annotate everything at once.
  \item \textbf{Context managers beyond \texttt{open()}.} \texttt{with pysam.AlignmentFile(\ldots) as bam:}, \texttt{with tempfile.TemporaryDirectory() as d:}, \texttt{with contextlib.suppress(FileNotFoundError):}. Any resource that needs deterministic cleanup --- file handles, temp dirs, database cursors, locks --- is a candidate.
\end{itemize}
\begin{tipbox}[Reach for the stdlib first]
Before adding a dependency, check \texttt{collections}, \texttt{itertools}, \texttt{functools}, \texttt{pathlib}, \texttt{dataclasses}, \texttt{contextlib}, \texttt{statistics}. The Python standard library solves most problems you'll hit in genomics scripting, and stdlib code ages well --- no surprise breaking changes when your analysis reviewer tries to rerun it two years later.
\end{tipbox}
\begin{winbox}[Checkpoint]
You wrote a function with a type-hinted signature, a dataclass for one biological record, and a \texttt{pathlib}-based file iteration. This is Python that will read well in a pull request.
\end{winbox}

\begin{tipbox}[The Python $\leftrightarrow$ Unix bridge]
You already know \texttt{awk '\{print \$1\}'} extracts column 1. The Python equivalent: \texttt{line.split('\textbackslash t')[0]}. Your Unix skills transfer directly. When in doubt, prototype in awk, then scale up in Python.
\end{tipbox}


\section{Block 4: Pandas, Plotting \& Bio Libraries}

\sprint{5}{Pandas for Tabular Genomic Data}{25}
\begin{itemize}[nosep]
  \item Load GWAS summary stats: \texttt{df = pd.read\_csv('sumstats.tsv', sep='\textbackslash t', dtype=\{'CHR': str\})}
  \item Filter: \texttt{sig = df[df['P'] < 5e-8]} --- genome-wide significant hits
  \item GroupBy: \texttt{df.groupby('CHR')['P'].min()} --- best p-value per chromosome
  \item Merge: \texttt{pd.merge(genotypes, phenotypes, on='SAMPLE\_ID')} --- the most common GWAS operation
\end{itemize}
\begin{winbox}[Checkpoint]
You loaded, filtered, and grouped real GWAS data. This replaces hours of spreadsheet work.
\end{winbox}

\sprint{6}{Visualization}{25}
\begin{itemize}[nosep]
  \item Manhattan plot: scatter plot of $-\log_{10}(P)$ vs.\ genomic position, colored by chromosome
  \item QQ plot: expected vs.\ observed $-\log_{10}(P)$ --- the standard GWAS diagnostic
  \item Use matplotlib for control, seaborn for aesthetics. Always label axes, always include titles.
  \item Save to file: \texttt{plt.savefig('manhattan.png', dpi=150, bbox\_inches='tight')}
\end{itemize}
\begin{winbox}[Checkpoint]
You made a Manhattan plot from real data. Screenshot it for your portfolio.
\end{winbox}

\sprint{7}{NumPy + SciPy = Fast Math}{25}
\begin{itemize}[nosep]
  \item NumPy arrays: 100$\times$ faster than Python lists for math. \texttt{np.mean(quality\_scores)} instead of \texttt{sum()/len()}
  \item Vectorized operations: never loop over millions of variants. \texttt{genotypes * 2} multiplies every element at once.
  \item SciPy stats: \texttt{scipy.stats.ttest\_ind(group1, group2)} for quick hypothesis tests
  \item Multiple testing: \texttt{from statsmodels.stats.multitest import multipletests} --- you'll need this for any -omics analysis
\end{itemize}
\begin{winbox}[Checkpoint]
You ran a statistical test on biological data in Python. No GUI, no spreadsheets --- just code.
\end{winbox}

\sprint{8}{Bio Libraries: pysam, Biopython, scanpy}{25}
\begin{itemize}[nosep]
  \item pysam: read BAM files programmatically. Count reads mapping to a gene in 3 lines.
  \item Biopython: \texttt{SeqIO.parse()} handles FASTA, FASTQ, GenBank. \texttt{Entrez.efetch()} queries NCBI.
  \item scanpy: the single-cell analysis framework. We'll use it heavily in Blocks 9--10.
  \item scikit-allel: population genetics in Python --- allele frequencies, LD, PCA from VCF data
\end{itemize}
\begin{winbox}[Checkpoint]
You queried NCBI from Python and read a BAM file programmatically. You're automating research.
\end{winbox}

\begin{restbox}[End of Blocks 3--4]
You just absorbed a new language AND applied it to biology. Take some time off. Let it consolidate.
\end{restbox}

\subsection{Checkpoint:}
\begin{itemize}[nosep]
  \item You can parse any common genomic file in Python (FASTA, BED, VCF, TSV)
  \item You can use Pandas to load, filter, group, and merge tabular data
  \item You can make a Manhattan plot and a QQ plot
  \item You understand why vectorized NumPy operations are faster than Python loops
\end{itemize}

\subsection{Resources}
\begin{itemize}[nosep]
  \item McKinney -- \textit{Python for Data Analysis} (the Pandas bible)
  \item VanderPlas -- \textit{Python Data Science Handbook} (free online)
  \item Real Python (realpython.com) --- searchable tutorials for anything
\end{itemize}
